\documentclass[10pt]{article}
\usepackage[T1]{fontenc}
\usepackage[utf8]{inputenc}
\usepackage[french]{babel}
\usepackage{amsmath,amssymb,mathrsfs,mathpazo}
\usepackage{geometry}
\geometry{a4paper,margin=2cm}
\usepackage{graphicx,array,tabularx,tabularray,colortbl,multirow,multicol,booktabs,enumitem,pifont,ulem,fancybox,diagbox,cancel}
\usepackage{pstricks,pstricks-add,pst-plot,pst-tree,pst-node,pst-text,pst-func,pst-math,pst-bspline,pst-3dplot,pst-eucl,pst-solides3d}
\let\pscalculate\undefined
\usepackage{pst-bezier}
\usepackage{tikz,pgfplots,tkz-tab}
\usetikzlibrary{arrows,arrows.meta,patterns,calc,positioning,trees,decorations.markings,shapes.geometric,shapes.misc,shapes.arrows,intersections,decorations.pathmorphing,decorations.pathreplacing,backgrounds,shadows,babel}
\pgfplotsset{compat=1.18}
\usepgfplotslibrary{fillbetween}
\usepackage[np]{numprint}
\usepackage{listings,makecell,pas-tableur,twoopt}
\let\barre\undefined
\usepackage[pstricks]{tablvar}
\definecolor{darkgreen}{RGB}{0,100,0}
\providecommand{\red}{\color{red}}
\providecommand{\blue}{\color{blue}}
\providecommand{\green}{\color{green}}
\providecommand{\black}{\color{black}}
\providecommand{\white}{\color{white}}
\providecommand{\gray}{\color{gray}}
\providecommand{\cyan}{\color{cyan}}
\providecommand{\magenta}{\color{magenta}}
\providecommand{\yellow}{\color{yellow}}
\providecommand{\violet}{\color{violet}}
\providecommand{\orange}{\color{orange}}
\providecommand{\euro}{\texteuro}
\providecommand{\EUR}{\texteuro}
\providecommand{\ieme}{\textsuperscript{e}}
\providecommand{\degres}{\ensuremath{^\circ}}
\providecommand{\degree}{\ensuremath{^\circ}}
\providecommand{\vv}[1]{\overrightarrow{#1}}
\providecommand{\Coord}[1]{#1}
\providecommand{\decosix}{}
\providecommand{\starredbullet}{\ensuremath{\bullet}}
\pagestyle{empty}
\setlength{\parindent}{0pt}
\usepackage{adjustbox}\DeclareUnicodeCharacter{2236}{:}\DeclareUnicodeCharacter{2190}{\ensuremath{\leftarrow}}
\providecommand{\R}{\mathbb{R}}
\providecommand{\N}{\mathbb{N}}
\providecommand{\Z}{\mathbb{Z}}
\providecommand{\Q}{\mathbb{Q}}
\providecommand{\D}{\mathbb{D}}
\providecommand{\C}{\mathbb{C}}
\providecommand{\e}{{\mathrm{e}}}
\providecommand{\vect}[1]{\overrightarrow{#1}}
\providecommand{\vectt}[1]{\overrightarrow{\mathrm{#1}}}
\providecommand{\vvt}[1]{\overrightarrow{\mathrm{#1}}}
\providecommand{\np}[2][]{{#2~#1}}
\providecommand{\vv}[1]{{\overrightarrow{#1}}}
\renewcommand{\d}{{\mathrm{d}}}
\providecommand{\degres}{{^\circ}}
\providecommand{\degre}{{^\circ}}
\providecommand{\blue}{\color{blue}}
\providecommand{\red}{\color{red}}
\providecommand{\enskip}{\quad}
\providecommand{\llbracket}{\lbrack\!\lbrack}
\providecommand{\rrbracket}{\rbrack\!\rbrack}
\providecommand{\euro}{{\text{€}}}
\providecommand{\no}{{\text{n°}}}
\providecommand{\interval}[2]{{[#1;#2]}}
\providecommand{\intervFO}[2]{{[#1;#2[}}
\providecommand{\dotfill}{\ldots}

\providecommand{\ds}{\displaystyle}
\providecommand{\textding}[1]{#1}
\providecommand{\Oij}{\(\mathrm{O};\vec{\imath},\vec{\jmath}\)}
\providecommand{\Oijk}{\(\mathrm{O};\vec{\imath},\vec{\jmath},\vec{k}\)}
\providecommand{\Ouv}{\(\mathrm{O};\vec{u},\vec{v}\)}
\providecommand{\intervFF}[2]{\ensuremath{[#1;#2]}}
\renewcommand{\euro}{\ifmmode\text{\texteuro}\else\texteuro\fi}
\DeclareUnicodeCharacter{00B0}{\ifmmode^{\circ}\else\textdegree\fi}
\DeclareUnicodeCharacter{20AC}{\ifmmode\text{\texteuro}\else\texteuro\fi}
% Énoncés : Éducation nationale. Transcription et corrigés : APMEP, Denis Vergès et contributeurs. Découpage et classement : CoopMaths, Sébastien Lozano, Éric Elter et contributeurs. Adaptation de la présentation : Maths Explorer.
% CC BY-SA : https://coopmaths.fr/www/about/
% https://coopmaths.fr/alea/static/bac/2025/tex/bac_2025_05_sujet1_ameriquenord_1_cor.tex
\begin{document}
\noindent\textit{APMEP / CoopMaths. CC BY-SA. Présentation adaptée par Maths Explorer.}\par\medskip
\medskip

\textbf{Partie A}

	\medskip

	\begin{enumerate}
		\item Puisqu'on s'intéresse à une connexion au hasard, on est en situation d'équiprobabilité, et les proportions annoncées dans l'énoncé sont assimilables à des probabilités.

Cela donne l'arbre pondéré ci-dessous :

	
	\begin{center}
	
	\begin{tikzpicture}[xscale=1,yscale=1]

\tikzstyle{fleche}=[->,>=latex,thick]
\tikzstyle{noeud}=[fill=white,circle,inner sep=2pt]
\tikzstyle{feuille}=[fill=white,circle,inner sep=2pt]
\tikzstyle{etiquette}=[pos=0.6,fill=white, inner xsep=3pt, inner ysep=1.5pt]

\def\DistanceInterNiveaux{3}
\def\DistanceInterFeuilles{0.75}

\def\NiveauA{(0)*\DistanceInterNiveaux}
\def\NiveauB{(1)*\DistanceInterNiveaux}
\def\NiveauC{(2)*\DistanceInterNiveaux}
\def\InterFeuilles{(-1)*\DistanceInterFeuilles}


\node[noeud] (R) at ({\NiveauA},{(2.5)*\InterFeuilles}) {};
\node[noeud] (Ra) at ({\NiveauB},{(0.5)*\InterFeuilles}) {$A$};
\node[feuille] (Raa) at ({\NiveauC},{(0)*\InterFeuilles}) {$S$};
\node[feuille] (Rab) at ({\NiveauC},{(1)*\InterFeuilles}) {$\overline{S}$};
\node[noeud] (Rb) at ({\NiveauB},{(2.5)*\InterFeuilles}) {$B$};
\node[feuille] (Rba) at ({\NiveauC},{(2)*\InterFeuilles}) {$S$};
\node[feuille] (Rbb) at ({\NiveauC},{(3)*\InterFeuilles}) {$\overline{S}$};
\node[noeud] (Rc) at ({\NiveauB},{(4.5)*\InterFeuilles}) {$C$};
\node[feuille] (Rca) at ({\NiveauC},{(4)*\InterFeuilles}) {$S$};
\node[feuille] (Rcb) at ({\NiveauC},{(5)*\InterFeuilles}) {$\overline{S}$};

\draw[fleche] (R.east)--(Ra.west) node[etiquette] {$0,25$};
\draw[fleche] (Ra.east)--(Raa.west) node[etiquette] {$0,9$};
\draw[fleche] (Ra.east)--(Rab.west) node[etiquette] {$0,1$};
\draw[fleche] (R.east)--(Rb.west) node[etiquette] {$0,15$};
\draw[fleche] (Rb.east)--(Rba.west) node[etiquette] {$0,8$};
\draw[fleche] (Rb.east)--(Rbb.west) node[etiquette] {$0,2$};
\draw[fleche] (R.east)--(Rc.west) node[etiquette] {$0,6$};
\draw[fleche] (Rc.east)--(Rca.west) node[etiquette] {$0,85$};
\draw[fleche] (Rc.east)--(Rcb.west) node[etiquette] {$0,15$};
\end{tikzpicture}
\end{center}


-

















\item La probabilité que la connexion soit stable et passe par le serveur B est $P(S \cap B)$.

$P(S \cap B) = P(B) \times P_B(S) = 0,15 \times 0,8 = 0,12$.

\item De même :

$P\left(C \cap \overline{S}\right) = P(C) \times P_C\left(\overline{S}\right) = 0,6 \times 0,15 = 0,09$.

Cela signifie que 9~\% des connexions à distance de l'entreprise sont des connexions transitant via le serveur C et qui sont instables.

\item Les évènements $A$, $B$ et $C$ forment une partition de l'univers, donc, d'après la formule des probabilités totales, on a :

$\begin{aligned}
P(S) &= P(A) \times P_A(S)+ P(B) \times P_B(S) + P(C) \times P_C(S)\\
&= 0,25 \times 0,9 + 0,15 \times 0,8 + 0,6 \times 0,85 \\
&= 0,225 + 0,12 + 0,51\\
&= 0,855
\end{aligned}$

La probabilité de l'évènement $S$ est donc bien $P(S)=0,855$.

\item La probabilité demandée est $P_S(B)$.
Par définition, on a :

$P_S(B) = \dfrac{P(S\cap B)}{P(S)} = \dfrac{0,12}{0,855} = \dfrac{8}{57}\approx \np{0,1403}$.

La probabilité que la connexion ait transité par le serveur B, sachant qu'elle est stable est d'environ 0,140, au millième près.
\end{enumerate}


\textbf{Partie B}

\medskip

\begin{enumerate}
\item 
	\begin{enumerate}
		\item Les éléments suivants ne sont pas nécessaires, puisqu'on admet que $X$ suit une loi binomiale :

\begin{itemize}
\item Chaque connexion est vue comme une expérience aléatoire à deux issues : le succès \og la connexion est instable \fg{}, de probabilité $p$

$p = P\left(\overline{S}\right) = 1-0,855 = 0,145$;
\item cette épreuve est répétée $n = 50$ fois, de façon supposée identique et indépendante, puisque la constitution de l'échantillon est réputée assimilable à un tirage avec remise;
\item $X$ est la variable aléatoire qui compte le nombre de succès, c'est-à-dire le nombre de connexions instables, sur ces 50 répétitions.
\end{itemize}

Donc $X$ suit la loi binomiale de paramètres $n = 50$ et $p = 0,145$.

		\item La probabilité qu'au plus huit connexions soient instables est $P(X \leqslant 8)$.

Avec la calculatrice, on a : $P(X \leqslant 8) \approx \np{0,7044}$, 

Donc la probabilité qu'au plus huit connexions soient instables est  $0,704$ arrondi au millième.
	\end{enumerate}

\item
	\begin{enumerate}
		\item Puisque $X_n$ suit la loi binomiale de paramètres $n=50$ et $p = 0,145$, pour tout entier naturel $k$ inférieur ou égal à $n$, on aura :

$P(X_n = k) = \displaystyle \binom{n}{k} \times 0,145^k \times 0,855^{n-k}$.

L'évènement \og au moins une connexion de cet échantillon est instable\fg{} est l'évènement contraire de \og aucune connexion de cet échantillon n'est instable\fg{}, autrement dit $\left\{X = 0\strut \right\}$.

Ainsi : $p_n = 1 - P(X_n = 0) = 1 - \displaystyle \binom{n}{0} \times 0,145^0 \times 0,855^{n-0} = 1 - 0,855^n$.

\item Résolvons, pour $n$ entier naturel :

$\begin{aligned}
p_n \geqslant 0,99
&\iff 1 - 0,855^{n} \geqslant 0,99\\
&\iff  - 0,855^{n} \geqslant -0,01\\
&\iff 0,855^{n} \leqslant 0,01\quad \text{car } -1<0\\
&\iff \ln\left(0,855^{n}\right) \leqslant \ln(0,01) \quad \text{car ln est croissante sur }\R^{*+}\\
	&\iff n\ln\left(0,855\right) \leqslant \ln(0,01)\\
	&\iff n \geqslant \dfrac{\ln(0,01)}{\ln\left(0,855\right)} \quad \text{car } 0,855 <1 \text{ donc }\ln\left(0,855\right) < 0\\
\end{aligned}$

Or $\dfrac{\ln(0,01)}{\ln\left(0,855\right)} \approx 29,4$, donc c'est à partir de $n = 30$ que l'on a une probabilité $p_n$ supérieure ou égale à 0,99.
\end{enumerate}

\item 
	\begin{enumerate}
		\item Puisque $X_n$ suit la loi binomiale de paramètres $n$ et $0,145$, alors l'espérance de $X_n$ est, par propriété : $E(X_n) = np = 0,145n$.

On a ensuite $F_n = \dfrac{1}{n} \times X_n$,

donc, par linéarité de l'espérance : $E(F_n) = \dfrac{1}{n} E(X_n) = \dfrac{1}{n}\times 0,145n = 0,145$.

On admet que $V\left(F_{n}\right)=\dfrac{\np{0,123975}}{n}$.

		\item Écrivons l'inégalité de concentration pour la moyenne empirique $F_n$, avec une précision $t = 0,1$ :

$\begin{aligned}
P\left(\left|F_{n}-E(F_n)\right| \geqslant t\right) \leqslant \dfrac{V(F_n)}{t^2}
	&\implies P\left(\left|F_{n}- 0,145\right| \geqslant 0,1\right) \leqslant \dfrac{~\dfrac{\np{0,123975}}{n}~}{0,1^2}\\
	&\implies P\left(\left|F_{n}-0,145\right| \geqslant 0,1\right) \leqslant \dfrac{\np{0,123975}}{n}\times \dfrac{1}{0,01}\\
	&\implies P\left(\left|F_{n}-0,145\right| \geqslant 0,1\right) \leqslant \dfrac{12,3975}{n}\\
				&\implies P\left(\left|F_{n}-0,145\right| \geqslant 0,1\right) \leqslant \dfrac{12,5}{n}\\
	& \qquad \text{car } \np{12,3975} < 12,5
\end{aligned}$

On a donc bien l'inégalité annoncée.

		\item Si on étudie un échantillon de \np{1000} connexions et que l'on constate que pour cet échantillon $F_{\np{1000}}=0,3$, alors on a $0,3 - 0,145 = 0,155$ : cette conduite de l'expérience réalise l'évènement $\Big\{ \left|F_{\np{1000}} - 0,145\strut \right| \geqslant 0,1\Big\}$.

D'après la question précédente, la probabilité de cet évènement est inférieure ou égale à $\dfrac{12,5}{\np{1000}} = \np{0,0125}$, il est donc hautement probable que les serveurs dysfonctionnent, car si la modélisation (basée sur un fonctionnement normal des serveurs) est fiable, une telle fréquence est très peu probable.
	\end{enumerate}
\end{enumerate}

\bigskip

\end{document}
